\centerline{\bf \Huge ДЗ$\cdot$Инварианты}
\title{ДЗ$\cdot$Инварианты}

\begin{flushright}
    \scriptsize
    $-$Анджур...Ты любишь монетки, Анджур?\\
    $-$НеНачальник...Я ИХ БЛИН НЕНАВИЖУ!\\
    Имя стенда $-$ Elistar Platinum. Имя мастера $-$ Анджур Аркадьевич. 
    
\end{flushright}


\problem{1}В компании ЛМШ работают 2022 лосося и 2023 шиншиллы в компанию пришли новые 2024 сотрудника. Руководитель компании распределяет новых сотрудников либо к шиншиллам либо к лососям. Может ли случиться так, что лососей и шиншилл будет поровну?

\problem{2}Изобретатели Анджур Аркадьевич и Расул Владимирович создали два автомата по изменению карточек и установили их в Перспективе. Первый, принимая карточку $(a; \ b; \ c; \ d; \ e)$, выдавал карточку $\biggl( \dfrac{a+b}{4}; \ \dfrac{b+c}{4}; \ \dfrac{c+d}{4}; \ \dfrac{d+e}{4}; \ \dfrac{e+a}{4} \biggl)$, а второй выдавал карточку ($a+b+c; \ b+c+d; \ c+d+e; \ d+e+a; \ e+a+b$). Во время очных сборов ЛМШ один из НеУчеников решил, используя эти два автомата, попробовать получить из карточки $(1; \ 2; \ 3; \ 4; \ 5)$ карточку $(25; \ 3; \ 2002; \ 16; \ 7)$. Докажите, что как бы не пытался НеУченик, у него не получится этого сделать.

\problem{3}Анджур Аркадьевич решил отомостить НеНачальнику за то, что он заскамил его учеников и пробрался в его логово. Конечно же злодей хорошо знал об этом и оставил для Анджура Аркадьевича испытание. Перед ним стояли огромные, бездонные весы, а рядом с ними два автомата. Первый принимал настоящую монету, а отдавал две фальшивые. Второй наоборот $-$ принимал фальшивую монету и выдавал две настоящие. Чтобы пройти дальше, нужно положить на левую чашу только настоящие монеты, а на правую только фальшивые таким образом, что весы встанут в равновесие. Стенд НеНачальника <<Фальшивомонетчик>> настолько силён, что даже его фальшивые монеты не отличаются по весу от настоящих. Сможет ли Анджур Аркадьевич пройти испытание, если изначально он пришёл с одной монеткой в кармане.